\subsection{Symplectic Vector Spaces}

\begin{lemma}
    Let $(V,\omega)$ be a symplectitc vector space and $W \subseteq V$ an isotropic subspace. Then any basis of $W$ can be extended to a symplectic basis of $V$.
\end{lemma}
\begin{proof}
    We proof this by induction over dim$(V)$. If dim$(V) = 2$ the statement follows directly from non-degeneracy of $\omega$. Now assume the statement holds for dim$(V) = 2(n-1)$ and let dim$(V) = 2n$. Given a basis $u_1,\ldots,u_k$ of $W$ define 
    \[
    U:= \mathrm{span}(u_2,\ldots,u_k) \subsetneq W
    \]
    Then $W^\omega \subsetneq U^\omega$, hence there exists $v_1 \in U^\omega$ with $v_1 \notin W^\omega$, so by construction we must have 
    \[
    \omega(u_1,v_1) \neq 0.
    \]
    After normalizing we may assume $\omega(u_1,v_1) = 1$. Now consider 
    \[
    Q := \mathrm{span}(u_1,v_1).
    \]
    Then $Q$ is a symplectic subspace, thus 
    \[
    (Q^\omega, \sigma := \omega|_{Q^\omega})
    \]
    again a symplectic vector space. Moreover $U \subseteq Q^\omega$ and
    \[
    U \subseteq U^\sigma.
    \]
    So by our induction hypothesis we can extend $u_2, \ldots, u_k$ to a symplectic basis $u_2,\ldots, u_n, v_2,\ldots,v_n$ of $(Q^\omega, \sigma)$. Then 
    \[
    u_1,\ldots,u_n,v_1,\ldots,v_n
    \] 
    is a symplectic basis of $(V,\omega)$.
\end{proof}

The following Lemma gives a useful characterisation of symplectic matrices analogous to orthogonal matrices.

\begin{lemma}
    A matrix $A \in \R^{2n \times 2n}$ is symplectic if and only if the column vectors of $A$ form a symplectic basis of $(\R^{2n}, \omega_0)$.
\end{lemma}
\begin{proof}
    Write $A = (a_{x_1} \cdots a_{x_n} a_{y_1} \cdots a_{y_n})$ for $a_{x_i}, a_{y_i} \in \R^{2n}$. If $A$ is symplectic we have 
    \[
    \omega_0(a_{x_i}, a_{y_j}) = \omega_0(Ae_{x_i},Ae_{y_i}) = \omega_0(e_{x_i},e_{y_i}) = \delta_{ij}
    \]
    and similiarly $\omega_0(a_{x_i}, a_{x_j}) = 0 = \omega_0(a_{y_i}, a_{y_j})$. For the other implication it suffices to show that $\omega_0(\cdot,\cdot) = \omega_0(A\cdot,A\cdot)$ holds on basis vectors, which follows from the same calculation.
\end{proof}